\scanpage{74}
Supposons maintenant la condition suivante satisfaite (en plus de : \uline{tout autom de} $E$ \uline{est intérieur}).

(C) Pour tout $o\subset E$, $o\in O$, la loi de composition de $\mathrm{Hom}_{\underline B_1}(o,o)=$ ensemble des \inserted{$(o,o)$-}biréseaux spéciaux \deleted{\unsure{sur $(o,o)$}} \deleted{est} $(L,M\longmapsto L.M)$ est commutative :

\[
LM=ML.
\]

% slanted
\margin{Il suffit (et il faut si $A$ est \unsure{Dedek.} de dim 1) que ce soit vrai sur les localisés \deleted{\unsure{en les}} idéaux maximaux (ou leurs complétés, si $A$ noeth.). Dans le cas local, \deleted{\unsure{il suffit}} c'est équiv. : les relations $LM=ML$ sont vraies pour $L,M$ des idéaux bilatères biprincipaux de $o$}

$C\in\mathrm{Ob}_0\,\underline{\underline B}_0$ une catégorie de Deuring, et $X\in\mathrm{ob}\,C$, de sorte que $o\unsure{\simeq|=}\mathrm{End}_C(X)\in O$ ; on peut \deleted{supposer} \inserted{choisir} $E\simeq o_K$ \inserted{(et ``identifier'')}. On va définir sur la gr-catégorie $\underline{\unsure{\mathrm{EquAut}|\mathrm{Equ}|\mathrm{Eqv}}}(C)=\underline{\mathrm{Hom}}_{\underline{\underline B}_0}(C,C)$ une \deleted{loi de} contrainte de commutativité, \deleted{\unsure{définit}} \inserted{en fait une \uline{cat. de Picard stricte}} à priori du choix de $X$. En effet, on a un gr-isomorphisme

\[
\underline{\mathrm{Hom}}_{\underline{\underline B}_0}(C,C)
\overset{\alpha_X}{\simeq}
\underline{\mathrm{Hom}}_{\underline{\underline B}_1}(C,C)
=\text{catégorie des }\deleted{\text{\unsure{bimodules} spéciaux sous }(o,o)}
\inserted{\text{$(o,o)$-réseaux}}\text{ spéciaux}
\]

cette dernière catégorie étant \inserted{cat. de Picard} strictmt comm. (en fait, les contraintes d'associativité \uline{et} de commutativité y sont des identités). On peut donc définir une contrainte de comm. stricte sur $\underline{\mathrm{Equ}}(C)$ par transport de structure. On peut l'expliciter ainsi : si $L,M$ sont des $(o,o)$-bimodules spéciaux, on définit

\[
c:L\otimes_o M\overset{\sim}{\longrightarrow}M\otimes_o L
\]

en \deleted{\unsure{défini}} utilisant des isom de $(o,o)$ bimodules

\[
\begin{aligned}L&\simeq L'\subset E\\M&\simeq M'\subset E\end{aligned}
\]

avec $L',M'$ des $(o,o)$-réseaux, et notant que

\[
\begin{aligned}
L\otimes M&\simeq L'M'.\\
&\phantom{\simeq L'}\shortparallel\\[-.5ex]
M\otimes L&\simeq M'L'
\end{aligned}
\]
% archive sheet number: 73

